% !TEX root = ../../main.tex

\section{DeepSeek}

DeepSeek 是大语言模型相关技术在系统中的应用形式之一。
大语言模型在自然语言理解、文本生成和语义组织方面具有较强能力，DeepSeek-V3 技术报告也体现了大语言模型在多类自然语言任务中的通用处理能力\cite{deepseek-v3-2024}。
在 HearBridge 系统中，DeepSeek 不直接参与视频视觉识别，也不替代手势识别模型，而是作为识别结果的语义后处理工具，用于对模型输出的词级候选结果进行整理和表达修正。

具体而言，Python 手势识别服务负责完成视频帧分析、关键点提取和词级识别结果生成；
Spring Boot 服务端根据识别结果构造语义修正请求，调用 DeepSeek 接口生成更符合自然语言表达习惯的句子级结果。
这样的设计能够将视觉识别和语言表达整理分开处理：识别模型负责判断视频中出现的手势动作，大语言模型负责在识别候选结果的约束下优化最终展示文本。

需要注意的是，DeepSeek 在本系统中的作用是辅助整理识别结果，而不是独立完成手语翻译。
如果原始识别结果错误较多，单纯依赖大语言模型可能会引入不符合实际动作内容的文本。
因此，本文将 DeepSeek 定位为语义修正和结果自然化工具，用于改善识别结果的可读性和用户理解体验。
